\chapter{Results}

This chapter presents the experimental results obtained from the proposed EEG-based
system. The work has two parts and the results follow the same split. The first part deals
with motor imagery, where the aim is to find out whether EEG recorded from a four-channel
dry headset carries enough information to tell a left-hand imagination apart from a
right-hand one. The second part deals with cognitive assessment, where the aim is to
compute attention, cognitive load, relaxation, stress and fatigue from the frequency bands
of the same signal. Results are reported in the order the work was carried out, starting
with the reproduction of the reference dataset analysis, then the analysis of the data
recorded in this project, then the classifier, and finally the cognitive scoring system.

\section{Evaluation Aspects}

For the motor imagery part, the EEG signals recorded during imagined left-hand and
right-hand grasps are compared in the time domain, through event-related desynchronisation
and synchronisation (ERD/ERS), and through power spectral density (PSD). These analyses
check whether the two classes show any consistent difference that a classifier can use. For
the cognitive part, the relative power in each frequency band is tracked over a session and
converted into scores, and the stability of the recording itself is checked through the
electrode impedance.

\subsection{Cases of Comparison}

The following comparison cases were studied during the experimental analysis:

\begin{enumerate}
    \item ERD/ERS of the reference dataset, reproduced to validate the analysis method
    \item ERD/ERS of the data recorded in this project, compared against the reference
    \item PSD of C3 and C4 for imagined left versus imagined right
    \item Classification accuracy of different SVM feature sets and kernels
    \item Band power and cognitive scores recorded over a complete session
\end{enumerate}

\subsection{Parameters Under Study}

The main parameters considered in this analysis are:

\begin{enumerate}
    \item EEG channels: C3, C4, A1 and A2, with C3 and C4 as the motor cortex pair
    \item Motor imagery classes: imagined left-hand grasp and imagined right-hand grasp
    \item Frequency range: 8--30 Hz for motor imagery, 1--30 Hz for cognitive scoring
    \item Features: band power, C3/C4 power ratio, mean, median, skewness and kurtosis
    \item Cognitive outputs: attention, cognitive load, relaxation, stress and fatigue
\end{enumerate}

\section{Experimental Setup}

The experiment was carried out using the GELO BCI headset with four EEG channels and one
reference electrode, sampling at 250 Hz. The acquisition was done through the Capsule GUI,
shown in Figure~\ref{fig:gui_acquisition}. Before every session the headset was paired over
Bluetooth, the electrode resistances were checked, and the two calibration steps described
in Chapter 5 were performed.

\begin{figure}[!htbp]
    \centering
    \includegraphics[width=0.85\linewidth]{gui_acquisition.jpg}
    \caption{EEG signal acquisition using the Capsule GUI}
    \label{fig:gui_acquisition}
\end{figure}

During each recording session the subject was seated comfortably and asked to imagine
gripping an object with the left hand or the right hand when prompted. The movement was
only imagined, so no actual movement was made. Each imagination was marked in real time
using a marker, which was recorded alongside the EEG so that the start and end of every
trial could be identified later. Each recorded file contains five marked repetitions of one
class. The session was saved in \texttt{.h5} format and converted to CSV for processing.

For post-recording analysis the CSV files were loaded and sorted by class based on their
filenames. The marker column was used to extract the active imagination windows as
individual trials. Since the shortest trial across the recordings was close to four
seconds, a fixed window of four seconds from the onset of each imagination was used so that
every trial contributes the same number of samples. The signal supplied by the device is
already filtered, so no further filtering was applied before the feature calculation except
the 8--30 Hz band limiting used for the ERD and PSD analysis. PSD was estimated using
Welch's method \cite{welch1967} with a window size of 256 samples at $F_s = 250$ Hz. Out of all the
recordings made, 50 right-hand trials and 22 left-hand trials were found usable, the rest
being discarded because of poor electrode contact or movement during the trial.

\FloatBarrier

\section{Observed Results}

\subsection{Reproduction of the Reference Dataset Analysis}

Before working on the recorded data, the ERD/ERS analysis of the reference
dataset~\cite{liu2024mi} was reproduced. This was done to confirm that the analysis method
used in this project gives the same result as the published work before it is applied to
new data. Event-related desynchronisation refers to a drop in EEG power in a given
frequency band during a task, and event-related synchronisation refers to a rise
\cite{pfurtscheller1999erd}. For motor
imagery, power in the 8--30 Hz range is expected to fall over the hemisphere opposite to
the imagined hand.

The C3 and C4 signals were extracted, filtered between 8 and 30 Hz using a second-order
zero-phase Butterworth filter, squared to obtain power, averaged across trials, and
smoothed. The change was computed as

\begin{equation}
\mathrm{ERD/ERS}~(\%) = \frac{X - M}{M} \times 100
\end{equation}

where $X$ is the power at a given instant and $M$ is the average power during the one
second preparation period before the task. Figure~\ref{fig:erd_paper} shows the result
obtained for subject 8 of the dataset. During left-hand imagery the C4 trace separates
clearly from C3 and falls through the middle of the trial, which is the expected
contralateral response. During right-hand imagery the two traces stay much closer together.
This matches the behaviour reported in the dataset paper and confirms that the processing
chain is correct.

\begin{figure}[!htbp]
    \centering
    \includegraphics[width=0.45\linewidth]{erd_paper.png}
    \caption[ERD/ERS of C3 and C4 for subject 8 of the reference dataset]%
            {ERD/ERS of C3 and C4 for subject 8 of the reference dataset~\cite{liu2024mi},
             reproduced using the analysis method of this project}
    \label{fig:erd_paper}
\end{figure}

\FloatBarrier

\subsection{ERD/ERS of the Recorded Data}

The same analysis was then applied to the data recorded with the GELO headset.
Figure~\ref{fig:erd_ours} shows the result for one subject over the four second imagination
window.

\begin{figure}[!htbp]
    \centering
    \includegraphics[width=0.95\linewidth]{erd_ours.jpg}
    \caption{ERD/ERS of C3 and C4 for the data recorded in this project, 8--30 Hz,
             $F_s = 250$ Hz}
    \label{fig:erd_ours}
\end{figure}

Two things can be seen from this figure. The first is that the two classes do separate. For
left-hand imagery the C3 trace stays above C4 for most of the window, reaching 20 to 40 per
cent, while for right-hand imagery the order reverses and C4 stays above C3, with both
traces remaining within roughly $\pm 15$ per cent. The sign of the difference between the
two channels therefore flips between the two classes, which is the property a classifier
can make use of.

The second is that the separation does not take the same form as in the reference dataset.
In the reference recording the contralateral channel falls below the baseline, giving a
clear desynchronisation, whereas in the recorded data both channels mostly rise above the
baseline and the difference between them is what carries the class information. There are
two likely reasons for this. The baseline period in the recorded sessions is taken from the
seconds immediately before the marker, during which the subject is already preparing for
the task, so the reference power is not a true rest level. In addition, the four-channel dry
headset picks up more muscle and movement artefact than the 30-channel saline system used
for the reference dataset, and these artefacts raise the power in the 8--30 Hz band. The
result is still usable, but it means the absolute ERD values from this setup should not be
compared directly against published values.

\FloatBarrier

\subsection{Power Spectral Density for Left and Right Imagery}

The PSD of the two motor cortex channels was computed for each trial.
Figure~\ref{fig:psd_left} shows a left-hand trial and Figure~\ref{fig:psd_right} a
right-hand trial from the same subject. In both figures the upper two plots show the two
channels separately and the lower plot overlays them, with the spectral peak marked.

\begin{figure}[!htbp]
    \centering
    \includegraphics[width=0.85\linewidth]{psd_left_mi.png}
    \caption{Power spectral density of C3 and C4 for an imagined left-hand grasp}
    \label{fig:psd_left}
\end{figure}

\begin{figure}[!htbp]
    \centering
    \includegraphics[width=0.85\linewidth]{psd_right_mi.png}
    \caption{Power spectral density of C3 and C4 for an imagined right-hand grasp}
    \label{fig:psd_right}
\end{figure}

\FloatBarrier

The difference between the two cases is clearest in the overlay plots. For the left-hand
trial the C3 curve sits roughly 10 dB above the C4 curve across most of the low frequency
range, giving a C3 to C4 power ratio of 9.52. For the right-hand trial the two curves lie
almost on top of each other, with a ratio of 0.86. The spectral peaks are close together in
both cases, at 5.50 Hz and 7.00 Hz for the left trial and 6.00 Hz and 6.50 Hz for the right
trial, which shows that the peak frequency itself does not separate the classes. It is the
balance of power between the two hemispheres, and not the shape of the spectrum, that
carries the useful information. This is the reason the logarithm of the C3 to C4 power
ratio was included as a feature for the classifier.

Both figures also show that the signal falls away sharply beyond about 45 Hz. This is the
effect of the filtering applied inside the device, and it is why the analysis in this work
is restricted to the 8--30 Hz range for motor imagery.

\FloatBarrier

\subsection{Classification Results}

Support Vector Machine classifiers were trained using the reference
dataset~\cite{liu2024mi} as training data and the data recorded in this project as test
data. Three feature sets were tried, each with both a linear and a radial basis function
kernel, and each using absolute power as well as ERD power:

\begin{itemize}
    \item \textbf{SVM-A:} mean, median, skewness and kurtosis of C3 and C4 (8 features)
    \item \textbf{SVM-B:} the above, plus $\log_{10}(P_{C3}/P_{C4})$ and the Welch band
          power over 8--30 Hz (9 features)
    \item \textbf{SVM-C:} the C3 to C4 power ratio alone (1 feature)
\end{itemize}

The accuracies obtained are summarised in Table~\ref{tab:svm_accuracy}.

\begin{table}[h]
\centering
\caption{Classification accuracy of the trained SVM models on the recorded data}
\label{tab:svm_accuracy}
\begin{tabular}{|l|l|c|c|}
\hline
\textbf{Model} & \textbf{Features} & \textbf{Absolute power} & \textbf{ERD power} \\
\hline
SVM-A & Mean, median, skewness, kurtosis of C3 and C4 & 53.9\% & 53.1\% \\
\hline
SVM-B & SVM-A features $+$ C3/C4 PSD ratio & 53.7\% & 53.5\% \\
\hline
SVM-C & C3/C4 PSD ratio only & 50.2\% & 51.4\% \\
\hline
\end{tabular}
\end{table}

All six combinations give accuracies between 50 and 54 per cent, which is close to what
would be obtained by guessing on a two-class problem. Adding the power ratio to the
statistical features in SVM-B did not improve the result, and using the ratio on its own in
SVM-C performed worst. Using ERD power instead of absolute power made very little
difference either way.

This result should be read together with the benchmarks reported for the reference dataset
itself, where CSP with LDA reached 55.6 per cent, filter bank CSP with an SVM reached 57.6
per cent and the best reported method reached 72.21 per cent, all using 30 channels on data
from the same recording system~\cite{liu2024mi}. Against that background, the figures in
Table~\ref{tab:svm_accuracy} are not surprising. Two difficulties combine here. The first
is that the model is trained on one recording system and tested on another, so the feature
distributions do not match: the training data comes from stroke patients recorded with 30
saline electrodes at 500 Hz, while the test data comes from healthy subjects recorded with
four dry electrodes at 250 Hz. The second is that only two channels are available for the
spatial information, so spatial filtering methods such as CSP \cite{ramoser2000csp}, which are what give the
published methods most of their accuracy, cannot be applied at all. The earlier ERD and PSD
results show that a difference between the classes does exist in the recorded data, so the
limitation lies in the training strategy and the number of channels rather than in the data
being empty of information.

\FloatBarrier

\section{Cognitive Assessment Results}

The second part of the system computes a cognitive score from the same recording. Results
for this part are reported from a complete monitoring session.

\subsection{Signal Quality and Filtering}

Since the whole cognitive score depends on band power, the quality of the contact between
the electrodes and the skin matters a great deal. Figure~\ref{fig:impedance} shows how the
resistance of the four electrodes changed over the first six minutes after the headset was
put on.

\begin{figure}[!htbp]
    \centering
    \includegraphics[width=0.8\linewidth]{electrode_impedance.png}
    \caption{Electrode resistance of the four channels after connecting the headset}
    \label{fig:impedance}
\end{figure}

All four electrodes start between 250 and 500 k$\Omega$ and settle to between 70 and 110
k$\Omega$ within about four minutes, which is well inside the usable range. C4 settles
highest and takes longest, which is consistent with it being the channel that gave the
noisiest motor imagery traces. The figure also shows why the session is not started
immediately: calibration is performed at around one minute and the session begins at about
one and a half minutes, by which time the contact has largely stabilised.

Two filters are applied to the continuous signal before the band powers are computed. Their
responses are shown in Figure~\ref{fig:filter}. A 0.5 Hz high-pass filter removes the slow
drift caused by sweat and small movements of the electrodes while leaving the delta band
intact, and a 50 Hz notch filter removes mains interference. Both filters process the
signal block by block while keeping their internal state, so that no discontinuity is
introduced at the block boundaries.

\begin{figure}[!htbp]
    \centering
    \includegraphics[width=0.8\linewidth]{filter_response.png}
    \caption{Frequency response of the 0.5 Hz high-pass and 50 Hz notch filters}
    \label{fig:filter}
\end{figure}

\FloatBarrier

\subsection{Usable Frequency Range}

Figure~\ref{fig:bandpower} shows how much of the original amplitude survives at each
frequency once the device filtering and the filters above are taken together.

\begin{figure}[!htbp]
    \centering
    \includegraphics[width=0.8\linewidth]{normalized_band_power.png}
    \caption{Fraction of the original amplitude retained across frequency bands}
    \label{fig:bandpower}
\end{figure}

The 7--13 Hz and 13--30 Hz ranges pass almost unchanged, at 1.00 and 0.96, and the 4--7 Hz
range is retained at 0.71. Below 4 Hz the response falls away sharply, to 0.14 in the 3--4
Hz range and effectively to zero below 2 Hz, and above 35 Hz it falls away in the same way.
The practical consequence is that the delta band cannot be measured reliably with this
device and the gamma band is only partly available, so the effective measurement range for
cognitive scoring is about 4 to 30 Hz. The scores used in this work are therefore built
from theta, alpha and beta, which all lie inside that range.

\subsection{Stress and Fatigue Reference Levels}

Stress is estimated from the arousal ratio, which is the relative high-beta power in the
20--30 Hz range divided by the relative alpha power \cite{gelo2026manual}. When a person
becomes tense, beta
activity rises and alpha activity falls, so this ratio increases.
Figure~\ref{fig:arousal} shows how the ratio behaved over one session.

\begin{figure}[!htbp]
    \centering
    \includegraphics[width=0.8\linewidth]{arousal_ratio.png}
    \caption{Arousal ratio over a session, with the rest and work reference levels marked}
    \label{fig:arousal}
\end{figure}

During the eyes-open baseline the ratio stays near 0.5, and a rest reference of 0.58 is
established from this period. Once work begins the ratio rises and moves between about 0.75
and 1.70 for the rest of the session. A work reference of 0.80 is taken from the first two
minutes of work, which are deliberately excluded from the stress calculation so that the
settling-in period is not reported as stress. The score itself is obtained by comparing the
current ratio against this reference,

\begin{equation}
\mathrm{Score} = 50 + 50 \tanh \left( \log_2 \frac{\mathrm{ratio}}{\mathrm{reference}} \right)
\end{equation}

which maps the ratio onto a 0 to 100 scale, with 50 meaning the person is at their own
reference level. The displayed value is the median over the last 30 seconds, so that a
single noisy block does not move the score. Fatigue is computed in the same way from the
ratio of combined theta and alpha power to beta power.

\FloatBarrier

\subsection{Complete Session Output}

Figure~\ref{fig:gui_cognitive} shows the interface at the end of a monitoring session, with
all the outputs computed from the single four-channel recording.

\begin{figure}[!htbp]
    \centering
    \includegraphics[width=\linewidth]{gui_cognitive_score.png}
    \caption{Cognitive score interface showing the session output}
    \label{fig:gui_cognitive}
\end{figure}

At the instant captured, attention is 91 and cognitive load is 87, both reported as high,
while relaxation is 23. Stress and fatigue are shown as not yet available, since both
require the opening period of the session to pass before a reference can be established.
The session plot in the middle shows how attention and cognitive load changed over the
minute of recording. The two do not move together throughout: around the half minute mark
cognitive load drops to near 10 while attention stays above 40, and both then climb back
together towards the end. This is useful because it shows the two measures are not simply
reporting the same quantity under different names.

The brain wave panel on the right shows the change in each band relative to the resting
baseline. Delta is down 92 per cent, theta down 83 per cent, alpha down 69 per cent, beta
down 50 per cent and gamma down 6 per cent. The pattern is consistent with an alert,
engaged state, since the slow bands associated with drowsiness are suppressed much more
than the faster bands associated with active thinking. The delta figure should not be read
too closely, as Figure~\ref{fig:bandpower} shows that very little of the delta band survives
the filtering on this device.

The sidebar also shows that the session follows a fixed sequence of setup, connection,
signal check, calibration, baseline and then the session itself, and the signal quality
indicator stayed green throughout, confirming that the contact held for the full recording.

\FloatBarrier

\section{Reason for the Observed Variations}

The differences seen between the two motor imagery classes come from the way the motor
cortex is organised. Imagining a movement of one hand changes the activity over the
opposite hemisphere, so the balance of power between C3 and C4 shifts depending on which
hand is imagined. This is what both the ERD traces in Figure~\ref{fig:erd_ours} and the
power ratios in Figures~\ref{fig:psd_left} and~\ref{fig:psd_right} show, and it is the only
property of the signal that separated the two classes reliably in this work.

The difference between the recorded data and the reference dataset has two main causes. The
headset used here has four dry electrodes while the reference system had thirty saline
electrodes, so both the signal quality and the amount of spatial information available are
lower. The baseline used for the ERD calculation is also taken from the moments just before
the marker, when the subject is already preparing for the task, so it does not represent a
true rest level and the computed percentages are shifted upward as a result.

The low classification accuracy follows from the same two causes, together with the decision
to train on one dataset and test on another. The subjects differ, the recording hardware
differs, the sampling rate differs and the number of channels differs, so features computed
on the training set do not fall in the same range as those computed on the test set. The
small amount of recorded data, 50 right-hand and 22 left-hand usable trials, also means the
two classes are unbalanced, which pushes a classifier towards the majority class and keeps
the accuracy near the chance level.

For the cognitive scoring, the variation between sessions is handled by referencing
everything to the person's own calibration. Because the scores are computed from relative
band powers and compared against a baseline recorded at the start of each session, a change
in overall signal level caused by electrode contact improving or worsening does not by
itself change the score.

\section{Summary of Analysis}

The results show that the EEG recorded from a four-channel dry headset does carry
information about which hand is being imagined. Both the ERD traces and the power spectral
density plots show a consistent difference between the two classes, with the C3 to C4 power
ratio reaching 9.52 for a left-hand trial and 0.86 for a right-hand trial on the same
subject. The direction of the difference between the two channels reverses between the
classes, which is the behaviour expected from the organisation of the motor cortex.

The classifiers trained so far do not convert this difference into useful accuracy. All six
combinations of feature set and kernel gave between 50 and 54 per cent on a two-class
problem. Since the published benchmarks on the reference dataset reach only 55 to 72 per
cent using thirty channels and matched training and test conditions, the result is in line
with the difficulty of the problem, but it confirms that training on a public dataset and
testing on data from a different device is not a workable strategy. The next step is to
record a larger and better balanced set of trials on the headset itself and to train and
test within that data, rather than across the two sources.

The cognitive assessment part of the system is working end to end. Electrode impedance
settles into the usable range within a few minutes, the filtering gives a reliable
measurement band of roughly 4 to 30 Hz, and attention, cognitive load, relaxation, stress
and fatigue are computed from the band powers and displayed live against the person's own
calibrated reference. The session output in Figure~\ref{fig:gui_cognitive} shows attention
and cognitive load varying independently over a recording, which indicates that the scores
are responding to the signal rather than tracking a single underlying quantity.

\section{Work carried out so far}

The work completed so far covers the following stages. A review of existing research was
carried out on EEG signal processing, motor imagery decoding, frequency band analysis and
machine learning methods for EEG, and the approach for the project was decided on the basis
of that review. The hardware was then set up using the GELO BCI headset, and a recording
procedure was established covering electrode impedance checking, alpha peak calibration,
baseline calibration and marker placement.

The ERD/ERS analysis reported in the reference dataset was reproduced to confirm that the
processing chain was correct, after which EEG data was recorded for left-hand and
right-hand motor imagery and the same analysis applied to it. The channels involved in left
and right motor imagery were identified, and the recorded data was compared against the
published results. Several SVM models were then trained using different feature sets and
kernels, and the features that separate the two classes were identified. In parallel, the
cognitive scoring system was built and tested over complete sessions, covering calibration,
band power computation, score calculation and the display interface.